\section{The operator}\label{cd:sec:operator}

\subsection{Semigroup and harmonic extension}

With $\widehat f(\omega)=\int f(x)e^{-i\omega x}dx$, the transform of $P_a$ is $e^{-a|\omega|}$, so
\[
P_a*P_b=P_{a+b},\qquad P_a=e^{-a|D|},\quad |D|=(-\partial_x^2)^{1/2}:
\]
the Cauchy densities form the Poisson semigroup, the symmetric stable semigroup of index one (for the clock of $\xi$ see Theorem~\ref{thm:poisson} and Proposition~\ref{prop:struct}, Section~\ref{ch:clock}). For an admissible measure $\mu$ the function
\[
h(x,y)=(P_y*\mu)(x),\qquad y>0,
\]
is harmonic in the upper half plane, being the real part of $\pi^{-1}i\int\bigl((z-t)^{-1}+t(1+t^2)^{-1}\bigr)d\mu(t)$ at $z=x+iy$. The integral converges absolutely because $(z-t)^{-1}+t(1+t^2)^{-1}=(1+tz)/((z-t)(1+t^2))=O_z(t^{-2})$, and the added term is real, so after multiplication by $i$ it changes only the imaginary part. (Without the added term the integral diverges for Lebesgue measure and for the zero counting measure, since $\sum_\gamma1/\lvert\gamma\rvert=\infty$.) The boundary values of $h$ are $\mu$ in the vague sense, so $P_a*\mu$ is the harmonic extension of $\mu$ restricted to the line $y=a$, and
\begin{equation}\label{cd:eq:deconv}
P_a*\mu=u
\end{equation}
asks for the boundary values of a harmonic function known on a parallel line, Hadamard's original example of an ill posed problem \cite{Hadamard23}. A second description, used in Chapter~\ref{chap:cx} (Proposition~\ref{cx:prop:subord}) and in Section~\ref{bp:sec:planar}, is probabilistic: $P_a$ is the law of planar Brownian motion started at height $a$ when it first hits the real axis, equivalently a Gaussian with a random variance distributed as the first passage time of one dimensional Brownian motion to level $a$.

\subsection{Uniqueness and the range}

\begin{proposition}[uniqueness]\label{cd:prop:unique}
If $\mu_1,\mu_2$ are admissible measures and $P_a*\mu_1=P_a*\mu_2$, then $\mu_1=\mu_2$.
\end{proposition}

\begin{proof}
Let $\nu=\mu_1-\mu_2$ and $g(x,y)=(P_y*\nu)(x)$, harmonic on $y>0$. For $y>a$, $g(\cdot,y)=P_{y-a}*g(\cdot,a)=0$; the semigroup property holds for the signed $\nu$ because $\lvert g(\cdot,a)\rvert\le P_a*\lvert\nu\rvert$ and $P_{y-a}*(P_a*\lvert\nu\rvert)=P_y*\lvert\nu\rvert<\infty$ by Tonelli, so Fubini applies. So $g$ vanishes on an open set and hence on the whole half plane. An admissible measure is the vague boundary value of its Poisson integral, so $\nu=0$.
\end{proof}

The inverse is formally multiplication by $e^{a|\omega|}$, which is not the Fourier transform of a tempered distribution, so the inverse is not a convolution and is defined only on the range of $P_a*$. If $u=P_a*\mu$, the Poisson extension of $u$ continues harmonically to the depth $a$ below the line, as $P_{y+a}*\mu$. The converse fails without conditions at depth $a$ (part (c) below). For positive sources Theorem~\ref{cd:thm:herglotz} characterises the range exactly. On $L^2(\R)$ the inverse is unbounded in every Sobolev norm, which is the source of everything that follows.

\begin{proposition}[the range]\label{cd:prop:range}
(a) For $f\in L^2(\R)$, $\int e^{2a|\omega|}|\widehat{P_a*f}|^2d\omega=\int\lvert\widehat f\rvert^2d\omega$, and every $u$ with $\int e^{2a|\omega|}|\widehat u|^2<\infty$ is $P_a*f$ for an $f\in L^2$. (b) For all real $s,t$, $\sup\|f\|_{H^t}/\|P_a*f\|_{H^s}=\infty$. (c) The function $u=2+\cos(x^2)$ is not $P_a*\mu$ for any $a>0$ and any admissible $\mu$, although its Poisson extension continues to an entire harmonic function.
\end{proposition}

\begin{proof}
(a) $\widehat{P_a*f}=e^{-a|\omega|}\widehat f$, and conversely $\widehat f=e^{a|\omega|}\widehat u$ defines $f\in L^2$ with $P_a*f=u$. (b) Let $f_N$ be real with $\widehat{f_N}$ the indicator of $[N,N+1]\cup[-N-1,-N]$. Then $\|P_a*f_N\|_{H^s}\le e^{-aN}\|f_N\|_{H^s}$, and $\|f_N\|_{H^t}/\|f_N\|_{H^s}$ lies between two powers of $N$, so $\|f_N\|_{H^t}/\|P_a*f_N\|_{H^s}\ge cN^{-|t-s|}e^{aN}\to\infty$. (c) The Poisson extension $h=P_y*u$ is continuous on $y\ge0$, and $g(x,y)=2+\Re\cos((x+iy)^2)$ is an entire harmonic function with $g(x,0)=u(x)$. So $h-g$ is harmonic on $y>0$, continuous on $y\ge0$ and zero on $y=0$; by the Schwarz reflection principle it extends to an entire harmonic function, and so does $h=g+(h-g)$. Suppose $u=P_a*\mu$ with $\mu$ admissible. Then $u'''=(\partial_x^3P_a)*\mu$ and $\lvert\partial_x^3P_a(x-t)\rvert\le C_a(1+(x-t)^2)^{-1}\le2C_a(1+x^2)(1+t^2)^{-1}$, because $1+t^2\le2(1+x^2)(1+(x-t)^2)$. So $\lvert u'''(x)\rvert\le C(1+x^2)$. But $u'''(x)=8x^3\sin(x^2)-12x\cos(x^2)$ is not $O(1+x^2)$.
\end{proof}

\section{Positivity}\label{cd:sec:positivity}

\subsection{The forward map preserves it, the inverse does not}

$P_a$ is a probability density, so $\mu\ge0$ implies $P_a*\mu\ge0$, and the semigroup property carries positivity up the levels: if $P_c*\mu\ge0$ for some $c$ then $P_{c'}*\mu\ge0$ for every $c'\ge c$. The converse fails, and the failure is quantitative.

\begin{example}[positive data, signed source]\label{cd:ex:signed}
Let $d>0$, $0<q<1$ and $\mu=\delta_0-q\delta_d$. Then $P_a*\mu=P_a-qP_a(\cdot-d)\ge0$ on $\R$ if and only if
\[
q\ \le\ \frac{r-d}{r+d},\qquad r=\sqrt{d^2+4a^2}.
\]
\end{example}

\begin{proof}
Positivity is $q\le\inf_xf(x)$ with $f(x)=(a^2+(x-d)^2)/(a^2+x^2)$. Setting $f'=0$ gives $x^2-dx-a^2=0$, and at the root $x^*=(d+r)/2$, using $x^{*2}=dx^*+a^2$, the numerator is $r(r-d)/2$ and the denominator $r(r+d)/2$. The other root is a maximum, and $f\to1$ at infinity.
\end{proof}

By Proposition~\ref{cd:prop:unique} the signed measure is the only source of these data at width $a$, so a positive datum can have no positive source at width $a$. The example has the structure of the Thorin family. Deconvolving $u=P_a*\mu$ only part of the way, by $P_c$ with $c<a$, gives $P_{a-c}*\mu=P_{a-c}-qP_{a-c}(\cdot-d)$, which is positive if and only if $q\le(r_c-d)/(r_c+d)$ with $r_c=\sqrt{d^2+4(a-c)^2}$; that is, if and only if
\[
c\ \le\ c^*=a-\tfrac12\sqrt{m^2-d^2},\qquad m=\frac{1+q}{1-q}\,d ,
\]
and $c^*\ge0$ if and only if the datum is positive: positivity holds up to a critical width and is lost beyond it. For $q=9-4\sqrt5\approx0.0557$, $a=\tfrac12$ and $d=1$ one has $m=\sqrt5/2$ and $c^*=\tfrac14$ exactly: the datum is positive (the bound of the example is $3-2\sqrt2\approx0.1716$), its partial deconvolutions are positive for $c<\tfrac14$, touch zero at $c=\tfrac14$ and are negative beyond. (Computed: with $q=0.0557$ the critical width is $0.25007$, so the value $\tfrac14$ needs $q=9-4\sqrt5$.)

\begin{proposition}[the critical width]\label{cd:prop:critical}
For $u\ge0$ with $\int u(x)(1+x^2)^{-1}dx<\infty$ let $c^*(u)=\sup\{c\ge0:\ u=P_c*\nu\text{ for some measure }\nu\ge0\}\in[0,\infty]$, where $P_0*\nu=\nu$, so that $c=0$ always qualifies. Then a positive source at width $c$ exists for every finite $c\le c^*(u)$ and for no $c>c^*(u)$; in particular the supremum is attained when it is finite.
\end{proposition}

\begin{proof}
If $u=P_c*\nu$ with $\nu\ge0$ and $0\le c'<c$, then $u=P_{c'}*(P_{c-c'}*\nu)$ with $P_{c-c'}*\nu\ge0$, by Tonelli and the semigroup property, so the admissible widths form an interval containing $0$. Let $0<c^*(u)<\infty$ and let $0<c_n\uparrow c^*(u)$ be admissible. Then $u=P_{c_1}*\nu_1$ is continuous, and by Theorem~\ref{cd:thm:herglotz} the harmonic extension $h$ of $u$ extends to a nonnegative harmonic function $V_n$ on $y>-c_n$. Since $V_n$ and $V_m$ agree with $h$ on $y>0$, they agree on the connected set $y>-\min(c_n,c_m)$, so together they define a nonnegative harmonic function on $y>-c^*(u)$ that extends $h$. Theorem~\ref{cd:thm:herglotz} with $c=c^*(u)$ gives $u=P_{c^*(u)}*\nu$ with $\nu\ge0$.
\end{proof}

For the Thorin family, let $u_1=\pi^{-1}\Re(\xi'/\xi)(1+iv)$, the density of level one. By Theorem~\ref{tf:thm:deconv}(a), with $\varepsilon_0=\tfrac12$ and $\varepsilon_1=\tfrac12-c$, a positive source at width $c\in(0,\tfrac12]$ exists if and only if $\zeta$ has no zeros in $\Re s>1-c$, zeros on $\Re s=1-c$ contributing atoms, and the source is the Thorin measure $\mathfrak U_{1-c}$, which is $u_{1-c}\,dv$ when it has no atoms. For $c>\tfrac12$ there is no source at all, since its Poisson integral would be harmonic on $y>-c$ and equal to $\pi^{-1}\Re(\xi'/\xi)(1+y+iv)$ off the zeros, which has poles at the zeros on the critical line (compare Theorems~\ref{lv:thm:thorin}, \ref{lv:thm:semigroup} and~\ref{lo:thm:dichotomy}). So
\[
c^*(u_1)=1-\sigma_{\max},\qquad\sigma_{\max}=\sup\{\Re\rho\},
\]
and the Riemann hypothesis is the statement $c^*(u_1)=\tfrac12$: the critical width of the level one density is as large as the functional equation allows; Remark~\ref{rem:ch2:oneline} says why positivity of the density on the single line is weaker.


\begin{proposition}[how much of the level one density a positive source explains]\label{cd:prop:share}
Let $\mathfrak u=u_1\,dv$, let $N(T)$ count the zeros with $0<\gamma\le T$ and $N(\sigma,T)$ those among them with $\beta>\sigma$. Then $\mathfrak u=P_{1/2}*\nu+\mathfrak d$ with $\nu\ge0$ and $\mathfrak d\ge0$, where, summing over all zeros with multiplicity,
\begin{gather*}
\nu=\sum_{\beta<1/2}P_{1/2-\beta}(\cdot-\gamma)+\sum_{\beta=1/2}\delta_\gamma+\sum_{\beta>1/2}2(1-\beta)\,\delta_\gamma,\\
\mathfrak d=\sum_{\beta>1/2}\bigl[P_{1-\beta}-2(1-\beta)P_{1/2}\bigr](\cdot-\gamma),
\end{gather*}
and
\begin{gather*}
\mathfrak d([-T,T])=4\int_{1/2}^1N(\sigma,T)\,d\sigma+O(\log^2T)=O(T),\\
\mathfrak u([-T,T])=2N(T)+O(\log^2T)\sim\frac T\pi\log T .
\end{gather*}
So the part of the mass of $\mathfrak u$ on $[-T,T]$ that is explained by a positive source at width $\tfrac12$ (a source $\nu\ge0$ with a nonnegative remainder) is $1-O(1/\log T)$, unconditionally. Status: proved.
\end{proposition}

\begin{proof}
A zero with $\beta<\tfrac12$ contributes $P_{1-\beta}=P_{1/2}*P_{1/2-\beta}$ to $\mathfrak u$, by the semigroup and $1-\beta>\tfrac12$. For $0<w<\tfrac12$, $P_w(x)-2wP_{1/2}(x)=\frac w\pi\bigl[(w^2+x^2)^{-1}-(\tfrac14+x^2)^{-1}\bigr]\ge0$ because $w^2<\tfrac14$, and with $w=1-\beta$ this treats a zero with $\beta>\tfrac12$, whose remainder has mass $1-2(1-\beta)=2\beta-1$. A zero on the line contributes $P_{1/2}*\delta_\gamma$. Summing the masses over the zeros with $\beta>\tfrac12$ and $|\gamma|\le T$ gives $4\sum(\beta-\tfrac12)=4\int_{1/2}^1N(\sigma,T)\,d\sigma$. Selberg's estimate \cite{Selberg46}, $N(\sigma,T)\ll T^{1-(\sigma-1/2)/4}\log T$, gives $\int_{1/2}^1N(\sigma,T)\,d\sigma\ll T\log T\int_0^{1/2}T^{-x/4}dx\le4T$. The terms $O(\log^2T)$ are the Cauchy tails of the zeros outside the window, and the count of $\mathfrak u([-T,T])$ is the Riemann von Mangoldt formula.
\end{proof}

The proposition carries no information about the Riemann hypothesis. The construction uses only a density estimate for the zeros to the right of the critical line, so the same share holds for any zero set with such an estimate, and a mass share does not separate $\zeta$ from functions whose off line zeros have density zero (heuristic for the Davenport Heilbronn function). The proportion of zeros known to lie on the critical line is likewise not the explained share. The Riemann hypothesis is the statement $\mathfrak d=0$ (Theorem~\ref{tf:thm:deconv}(b)). The proposition follows a referee's correction of an unrefereed note of Polson (November 2026), which had called a different decomposition the explained share.

\subsection{When does a positive source exist}

\begin{theorem}[existence of a positive source]\label{cd:thm:herglotz}
Let $u\ge0$ be continuous with $\int u(x)(1+x^2)^{-1}dx<\infty$, and let $h(x,y)=(P_y*u)(x)$ be its harmonic extension, $y>0$. For $c>0$ the following are equivalent.
\begin{enumerate}
\item[(i)] $u=P_c*\nu$ for a measure $\nu\ge0$.
\item[(ii)] $h$ extends to a nonnegative harmonic function on the half plane $y>-c$.
\end{enumerate}
\end{theorem}

\begin{proof}
If (i) holds, $V(x,y)=(P_{y+c}*\nu)(x)$ is harmonic and nonnegative on $y>-c$ and equals $h$ for $y\ge0$ by the semigroup property. Conversely, a nonnegative harmonic function $V$ on $y>-c$ is, by the Herglotz and Riesz representation for the half plane (see \cite{ABR01}), of the form $V(x,y)=\eta(y+c)+(P_{y+c}*\nu)(x)$ with $\eta\ge0$, $\nu\ge0$ and $\int(1+t^2)^{-1}d\nu<\infty$, and $\eta=0$: at $y=0$ the representation says $u=\eta c+P_c*\nu$, since $V=h$ on $y\ge0$ and $u$ is continuous, so for $y>0$ we have both $V(\cdot,y)=\eta(y+c)+P_{y+c}*\nu$ and $h(\cdot,y)=P_y*u=\eta c+P_{y+c}*\nu$, whence $\eta y=0$. Thus $u=P_c*\nu$.
\end{proof}

For the level one density the harmonic extension is explicit. Smoothing raises the level, $P_y*u_1=u_{1+y}$, and the formula continues below the line: $u_\theta$ is harmonic in $(\theta,v)$ wherever $\xi$ has no zeros. So condition (ii) for $u_1$ at width $c$ says that $\xi$ has no zeros on $\Re s>1-c$ and $\Re(\xi'/\xi)(s)\ge0$ there, and at $c=\tfrac12$ it is Lagarias's criterion \cite{Lagarias}: the Riemann hypothesis is equivalent to $\Re(\xi'/\xi)(s)>0$ for $\Re s>\tfrac12$ (the conditions $\ge0$ and $>0$ agree by the minimum principle).


\subsection{A Pick matrix test}\label{cd:sec:pick}

Positivity of the source at width $c$ says that $f=\xi'/\xi$ is a Carath\'eodory function (nonnegative real part) on a half plane, and Carath\'eodory functions are characterised by positive semidefinite Pick matrices (the Nevanlinna Pick theorem, \cite{Donoghue74}). This turns Theorem~\ref{cd:thm:herglotz} for the level one density into positivity of finite matrices built from $\xi'/\xi$ on the line $\Re s=1$. The statements of this subsection are from an unrefereed note of Polson (November 2026), with the proof shortened and the window form added after a referee's reading. No priority is claimed, and at $\lambda=\tfrac12$ the criterion is a restatement of Lagarias's \cite{Lagarias}.

For $\lambda\in[\tfrac12,1)$ and distinct nodes $z_j=1+i\tau_j$, $j=1,\dots,n$, put
\[
P^{(\lambda)}_{jk}=\frac{f(z_j)+\overline{f(z_k)}}{z_j+\overline{z_k}-2\lambda},\qquad f=\frac{\xi'}{\xi},
\]
the denominator being $2(1-\lambda)+i(\tau_j-\tau_k)\ne0$. On the line $\Re z=1$ the real part of $f$ is $\pi u_1>0$, so $f$ is determined there by the data up to an additive imaginary constant (Herglotz representation), which $\Im f(1)=0$ fixes because $\xi$ is real on $\R$. Thus $P^{(\lambda)}$ is computed from $u_1$ alone.

\begin{lemma}[Gram decomposition]\label{cd:lem:gram}
Let $k_\rho=\bigl((z_j-\rho)^{-1}\bigr)_j$, let $J$ be the all ones matrix, $S_{jk}=(z_j+\overline{z_k}-2\lambda)^{-1}$, and let $\circ$ be the entrywise product. Then, with the zeros summed in conjugate pairs,
\[
P^{(\lambda)}=\sum_\rho\bigl(k_\rho k_\rho^{*}\bigr)\circ\bigl(J+2(\lambda-\beta_\rho)S\bigr),
\]
each entry converging absolutely, and $S$ is positive semidefinite. Status: proved.
\end{lemma}

\begin{proof}
For one zero, $\bigl[(z-\rho)^{-1}+(\bar w-\bar\rho)^{-1}\bigr](z+\bar w-2\lambda)^{-1}=(z+\bar w-2\beta)\big/\bigl[(z-\rho)(\bar w-\bar\rho)(z+\bar w-2\lambda)\bigr]$, and $z+\bar w-2\beta=(z+\bar w-2\lambda)+2(\lambda-\beta)$. Each term is $O(\gamma^{-2})$, so the sum over zeros converges absolutely and agrees with the Hadamard sum for $f$. The matrix $S$ is positive semidefinite since $(z+\bar w-2\lambda)^{-1}=\int_0^\infty e^{-x(z-\lambda)}\,\overline{e^{-x(w-\lambda)}}\,dx$ for $\Re z,\Re w>\lambda$.
\end{proof}

\begin{theorem}[Pick criterion]\label{cd:thm:pick}
Let $\lambda\in[\tfrac12,1)$. The following are equivalent.
\begin{enumerate}
\item[(i)] $\xi$ has no zero with $\beta>\lambda$.
\item[(ii)] $u_1\,dv=P_{1-\lambda}*\nu$ for a measure $\nu\ge0$.
\item[(iii)] $P^{(\lambda)}\succeq0$ for every finite set of distinct nodes on the line $\Re z=1$.
\item[(iv)] $P^{(\lambda)}\succeq0$ for every finite set of distinct nodes in one segment $\{1+i\tau:\ a\le\tau\le b\}$, $a<b$.
\end{enumerate}
Status: proved.
\end{theorem}

\begin{proof}
(i)$\Leftrightarrow$(ii) is Theorem~\ref{tf:thm:deconv}(a) with $\varepsilon_0=\tfrac12$ and $\varepsilon_1=\lambda-\tfrac12$ (see Proposition~\ref{cd:prop:critical}). (i)$\Rightarrow$(iii): if $\beta_\rho\le\lambda$ for every zero, then $J+2(\lambda-\beta_\rho)S$ is positive semidefinite, so is its entrywise product with the rank one matrix $k_\rho k_\rho^*$ (Schur product theorem), and so is the sum, by Lemma~\ref{cd:lem:gram}. (iii)$\Rightarrow$(iv) is trivial. (iv)$\Rightarrow$(i): let $H_\lambda=\{\Re z>\lambda\}$ and let $z_1,z_2,\dots$ be distinct nodes, dense in the segment. By the Nevanlinna Pick theorem for the half plane \cite{Donoghue74}, for each $n$ there is a holomorphic $g_n$ on $H_\lambda$ with $\Re g_n\ge0$ and $g_n(z_j)=f(z_j)$ for $j\le n$. Since $g_n(z_1)=f(z_1)$ is fixed, the $g_n$ form a normal family (Carath\'eodory functions with a fixed value at one point, by Harnack's inequality). A subsequence converges locally uniformly to a holomorphic $g$ on $H_\lambda$ with $\Re g\ge0$ and $g(z_j)=f(z_j)$ for every $j$. Now $f$ is meromorphic on $\C$, has no pole on $\Re z=1$, and agrees with $g$ on a set with a limit point in $H_\lambda$, so $f=g$ on $H_\lambda$. Hence $f$ has no pole there, which is (i).
\end{proof}

\begin{corollary}[one sided test]\label{cd:cor:cert}
If $P^{(\lambda)}$ has a negative eigenvalue for some finite set of nodes, then $\xi$ has a zero with $\beta>\lambda$. Evaluated in ball arithmetic this would be a certified proof of that statement. Status: proved.
\end{corollary}

For $\xi$ itself and $\lambda\ge\tfrac12$ a negative eigenvalue would disprove the Riemann hypothesis, so the test is informative for model functions and for other $L$ functions. At $\lambda=\tfrac12$ under the Riemann hypothesis the Gram decomposition reduces to $P^{(1/2)}=\sum_\rho k_\rho k_\rho^*$, a sum of rank one matrices, so the matrix of the data is the spectral object of the zeros.

\begin{table}[ht]
\centering
\caption{Smallest eigenvalue of $P^{(\lambda)}$, computed in $60$ digit arithmetic (not ball arithmetic) at the $80$ nodes $1\pm i\tau_j$, with $40$ ordinates $\tau_j$ equally spaced in $[10,50]$. For the model, $f$ is the logarithmic derivative of the function with the first $30$ zeros on the line and a pair at $\beta=0.8$ and $0.2$ at height $35$. Status: computed.}\label{cd:tab:pick}
\begin{tabular}{lccc}
\toprule
$\lambda$ & $0.51$ & $0.50$ & $0.49$\\
\midrule
true $\xi$ & $2.60\times10^{-4}$ & $-4.8\times10^{-61}$ & $-0.0513$\\
\bottomrule
\end{tabular}

\medskip
\begin{tabular}{lccc}
\toprule
$\lambda$ & $0.62$ & $0.63$ & $0.65$\\
\midrule
model with a pair at height $35$ & $-0.0838$ & $-0.0382$ & none negative\\
\bottomrule
\end{tabular}
\end{table}

For the true $\xi$ the matrix has no negative pivot for $\lambda\ge0.501$, and at $\lambda=\tfrac12$ the smallest eigenvalue is zero to the working precision, which is a statement about the precision and not about the zeros (under the Riemann hypothesis $P^{(1/2)}$ is a Gram matrix, hence positive semidefinite and numerically singular). It is negative below $\tfrac12$, where the theorem does not apply and positivity fails because of the poles on the critical line. For the model the matrix has two negative eigenvalues for every $\lambda\in(\tfrac12,0.63]$ and none at $0.65$, so by Corollary~\ref{cd:cor:cert} the test detects a zero of the model with $\beta>0.63$ while the true value is $0.8$. The threshold depends on the nodes and moves towards $0.8$ as nodes are added near height $35$.

\section{Stability}\label{cd:sec:stability}

\subsection{H\"older between levels}

\begin{theorem}[three lines]\label{cd:thm:holder}
Let $b_0<b<a$ and let $\mu$ be a signed measure with $P_{b_0}*\mu\in L^2$, where for $b_0<0$ this means $\mu=P_{|b_0|}*\lambda$ with $\lambda\in L^2$, and $P_{b_0}*\mu$ stands for $\lambda$. Then
\[
\|P_b*\mu\|_2\ \le\ \|P_a*\mu\|_2^{\,\frac{b-b_0}{a-b_0}}\ \|P_{b_0}*\mu\|_2^{\,\frac{a-b}{a-b_0}} .
\]
In particular, if two sources agree in their data at width $a$ to within $\varepsilon$ in $L^2$ and both are bounded by $M$ at width $b_0$, they agree at width $b$ to within $\varepsilon^{\kappa}(2M)^{1-\kappa}$, $\kappa=(b-b_0)/(a-b_0)$. The exponent is attained in rate.
\end{theorem}

\begin{proof}
Let $F$ be the Fourier transform of $P_{b_0}*\mu\in L^2$ (of $\lambda$ when $b_0<0$). For $y>b_0$, $P_y*\mu=P_{y-b_0}*(P_{b_0}*\mu)$, so $\|P_y*\mu\|_2^2=(2\pi)^{-1}\int e^{-2(y-b_0)|\omega|}|F|^2d\omega$, and with $\kappa=(b-b_0)/(a-b_0)$,
\[
e^{-2(b-b_0)|\omega|}|F|^2=\bigl(e^{-2(a-b_0)|\omega|}|F|^2\bigr)^{\kappa}\bigl(|F|^2\bigr)^{1-\kappa},
\]
so H\"older's inequality with exponents $1/\kappa$ and $1/(1-\kappa)$ gives the bound. Equivalently, $y\mapsto\log\|P_y*\mu\|_2$ is convex, the three lines theorem for the harmonic extension. For sharpness take $F$ concentrated near a frequency $\Omega$ with $e^{-(a-b_0)\Omega}=\varepsilon/M$: both sides are then equal up to constants.
\end{proof}

For the Thorin family under the Riemann hypothesis the data are at level $1$, which is width $a=\tfrac12$ above the zero counting measure, and the level $\theta'$ measure is at width $b=\theta'-\tfrac12$. The zero measure is atomic, so Theorem~\ref{cd:thm:holder} admits no level below it ($b_0<0$ would require $\sum_\gamma\delta_\gamma=P_{\lvert b_0\rvert}*\lambda$ with $\lambda\in L^2$, a smooth function), and it has infinite mass. For the zeros in a window, a finite measure, Proposition~\ref{cd:prop:finite} below with $w=0$, $a=\tfrac12$ and $b=\theta'-\tfrac12$ gives the exponent $2\theta'-1$: about $\tfrac12$ at $\theta'=\tfrac34$, falling to $0$ at $\theta'=\tfrac12$. The application is heuristic, since a window does not commute with $P_y$ and the zeros outside it carry infinite mass; heuristically, it explains why $\widetilde U_{3/4}$ was recovered so much better than $\widetilde U_{1/2}$ in Chapters~\ref{chap:lo} and~\ref{chap:cx}. Table~\ref{cd:tab:slopes} shows the exponent $(b+w)/(a+w)$ of Proposition~\ref{cd:prop:finite} on a test source of Lorentzians of width $w$, to within $0.01$.

\begin{proposition}[finite sources]\label{cd:prop:finite}
Let $\nu$ be a finite signed measure of total variation at most $B$, let $w\ge0$ and $h=P_w*\nu$ (with $P_0*\nu=\nu$, and $P_b*h$ meaning $P_{b+w}*\nu$ when $b<0$), and let $-w<b<a$. Then
\[
\|P_b*h\|_2\le\sqrt2\,D^{1-\kappa}\,\|P_a*h\|_2^{\,\kappa},\qquad\kappa=\frac{b+w}{a+w},\quad D=\frac{B}{\sqrt{2\pi(b+w)}} ,
\]
and no larger exponent holds uniformly over this class.
\end{proposition}

\begin{proof}
The transform of $h$ satisfies $|\widehat h(\omega)|\le Be^{-w|\omega|}$. Let $\varepsilon=\|P_a*h\|_2$ and $\Omega\ge0$. For $|\omega|\le\Omega$, $e^{-2b|\omega|}=e^{2(a-b)|\omega|}e^{-2a|\omega|}\le e^{2(a-b)\Omega}e^{-2a|\omega|}$, so by Plancherel
\[
\begin{split}\|P_b*h\|_2^2&=\frac1{2\pi}\int e^{-2b|\omega|}|\widehat h(\omega)|^2d\omega\\&\le e^{2(a-b)\Omega}\varepsilon^2+\frac{B^2}{2\pi}\int_{|\omega|>\Omega}e^{-2(b+w)|\omega|}d\omega\\&=e^{2(a-b)\Omega}\varepsilon^2+D^2e^{-2(b+w)\Omega}.\end{split}
\]
If $\varepsilon<D$, take $e^{(a+w)\Omega}=D/\varepsilon$: both terms equal $D^{2(1-\kappa)}\varepsilon^{2\kappa}$. If $\varepsilon\ge D$, take $\Omega=0$: then $\|P_b*h\|_2\le D\le D^{1-\kappa}\varepsilon^\kappa$. For sharpness let $\phi(x)=(2\pi)^{-1}\bigl(\sin(x/2)/(x/2)\bigr)^2$, so that $\phi\ge0$, $\int\phi=1$ and $\widehat\phi(\omega)=(1-|\omega|)_+$, and let $\nu_\Omega=B\phi(x)\cos(\Omega x)\,dx$ with $\Omega\ge2$, of total variation at most $B$, and $h_\Omega=P_w*\nu_\Omega$. Then $\widehat{\nu_\Omega}=\tfrac B2\bigl(\widehat\phi(\cdot-\Omega)+\widehat\phi(\cdot+\Omega)\bigr)$ has two disjoint bumps, and for every $y>-w$, $\|P_y*h_\Omega\|_2=Be^{-(y+w)\Omega}c_\Omega(y)$ with $(6\pi)^{-1/2}e^{-(y+w)}\le c_\Omega(y)\le(6\pi)^{-1/2}e^{y+w}$. So $\|P_b*h_\Omega\|_2/\|P_a*h_\Omega\|_2^{\,\kappa}$ lies between two positive multiples of $B^{1-\kappa}$, independent of $\Omega$, while $\|P_a*h_\Omega\|_2\to0$ as $\Omega\to\infty$; this excludes every exponent above $\kappa$.
\end{proof}

With $w=0$ this bounds a finite signed measure from data at width $a$; with $b=0$ and $w>0$ it gives the exponent $w/(a+w)$ for recovering $h=P_w*\nu$ from the data at width $a$: the improvement comes from the analytic smoothing and the mass bound, not from positivity, and if $w\to0$ or the mass bound is removed no uniform positive exponent follows.


\begin{proposition}[the exponent is attained inside the positive cone]\label{cd:prop:conelower}
For $0<b<a$ there is a constant $k>0$ depending only on $a,b$ such that for every small $\delta>0$ there are positive measures $\nu_0,\nu_1$ of mass at most $2$ with $\|P_a*(\nu_1-\nu_0)\|_2\le\delta$ and $\|P_b*(\nu_1-\nu_0)\|_2\ge k\,\delta^{b/a}$. Status: proved.
\end{proposition}

\begin{proof}
Take $r=2a$, $\nu_0=P_r$ and $\nu_1=P_r(x)(1+\cos\Omega x)\ge0$. The difference has transform $\tfrac12\bigl(e^{-r|\omega-\Omega|}+e^{-r|\omega+\Omega|}\bigr)$, so for $s\in\{a,b\}$ and $\Omega\ge1$ the norm $\|P_s*(\nu_1-\nu_0)\|_2$ is comparable, with constants depending on $a,b$, to $e^{-s\Omega}$, because the frequencies near $\pm\Omega$ dominate when $r>s$. Choose $\Omega$ with $e^{-a\Omega}\asymp\delta$.
\end{proof}

So inside the positive cone the exponent $b/a$ of Proposition~\ref{cd:prop:finite} (with $w=0$) is attained: positivity costs nothing in the rate and gains nothing, the rate being a consequence of the mass bound (Corollary~\ref{cd:cor:positive} removes the need for a mass bound but not the exponent). The proposition is from an unrefereed note of Polson (November 2026), whose own stability theorem for positive measures is the weaker form of Proposition~\ref{cd:prop:finite}.

\begin{table}[ht]
\centering
\caption{Stability between levels, computed. The source is a sum of $40$ Lorentzians of width $w=0.05$ on a period $400$, the data are at width $a=\tfrac12$ with white noise of size $\varepsilon$ at every frequency, and $P_b*\mu$ is recovered by an oracle spectral cutoff, with $\varepsilon$ from $10^{-12}$ to $10^{-3}$. The fitted slope is that of $\log$ of the error in $P_b*\mu$ against $\log\varepsilon$; the predicted slope is the exponent $(b+w)/(a+w)$ of Proposition~\ref{cd:prop:finite}. Nineteen values of $\varepsilon$ (the error does not depend on the noise draw). All four fitted slopes are within $0.004$ of the prediction. Status: computed, illustration.}\label{cd:tab:slopes}
\begin{tabular}{lcccc}
\toprule
level $b$ & $\tfrac38$ & $\tfrac14$ & $\tfrac18$ & $0$ \\
\midrule
predicted $(b+w)/(a+w)$ & $0.7727$ & $0.5455$ & $0.3182$ & $0.0909$ \\
fitted slope & $0.769$ & $0.545$ & $0.318$ & $0.091$ \\
\bottomrule
\end{tabular}
\end{table}

If the datum and both sources are nonnegative, the mass in a window is controlled by the data alone, and the three lines theorem needs no a priori bound.

\begin{corollary}[positive sources]\label{cd:cor:positive}
If $\mu_1,\mu_2\ge0$, $u_j=P_a*\mu_j\in L^2$ and $\|u_1-u_2\|_2\le\varepsilon$, then for $0<d<b<a$
\[
\|P_b*(\mu_1-\mu_2)\|_2\le\varepsilon^\kappa\bigl((a/d)(\|u_1\|_2+\|u_2\|_2)\bigr)^{1-\kappa},\qquad\kappa=\frac{b-d}{a-d}.
\]
\end{corollary}

\begin{proof}
For $0<d<a$, $P_d/P_a=(d/a)(a^2+x^2)/(d^2+x^2)\le a/d$ pointwise, so $P_d*\mu_j\le(a/d)P_a*\mu_j$ pointwise, and $\|P_d*(\mu_1-\mu_2)\|_2\le(a/d)(\|u_1\|_2+\|u_2\|_2)$. Theorem~\ref{cd:thm:holder} applies with $b_0=d$.
\end{proof}

This needs data in $L^2$ on all of $\R$; for data on a window the question is open (Problem~\ref{cd:prob:three}).

\subsection{Logarithmic at the boundary}

\begin{theorem}[the boundary]\label{cd:thm:log}
Suppose $\|\mu\|_{H^s}\le M$ for some $s>0$ and $\|P_a*\mu\|_2\le\varepsilon$, with $\varepsilon/M$ small enough, depending on $a$ and $s$. Then
\[
\|\mu\|_2\ \le\ C_s\,M\Bigl(\frac{2a}{\log(M/\varepsilon)}\Bigr)^{s},
\]
and no bound of smaller order holds uniformly over the class: the modulus of continuity of the inverse is of this order.
\end{theorem}

\begin{proof}
Split the frequencies at $\Omega$. Below, $|\widehat\mu|\le e^{a\Omega}|\widehat{P_a*\mu}|$, contributing at most $e^{a\Omega}\varepsilon$ to the norm; above, the $H^s$ bound gives at most $M\Omega^{-s}$. With $\Omega=(2a)^{-1}\log(M/\varepsilon)$ the first term is $(M\varepsilon)^{1/2}$, of smaller order, and the second is the bound. For the lower bound, a source with transform concentrated near the frequency $\Omega=a^{-1}\log(M/\varepsilon)$ and $H^s$ norm $M$ has $L^2$ norm of order $M\Omega^{-s}$ and data of size at most $\varepsilon$.
\end{proof}

The sharp modulus $\sup\{\|\mu\|_2:\ \|\mu\|_{H^s}\le M,\ \|P_a*\mu\|_2\le\varepsilon\}$ equals $(1+o(1))\,M\bigl(a/\log(M/\varepsilon)\bigr)^s$ as $\varepsilon/M\to0$: in the upper bound take $\Omega=a^{-1}\bigl(L-(s+1)\log L\bigr)$ with $L=\log(M/\varepsilon)$, which makes the low frequency term $ML^{-s-1}$ and the high frequency term $(1+o(1))M(a/L)^s$; the lower bound is the one given.

The logarithm is sharp for positive measures of bounded mass as well. The next theorem is from an unrefereed note of Polson (November 2026). Write $d_{\rm BL}(\mu,\mu')=\sup\{\int\varphi\,d(\mu-\mu'):\ \|\varphi\|_\infty\le1,\ \mathrm{Lip}(\varphi)\le1\}$ for the bounded Lipschitz distance.

\begin{theorem}[Sharp logarithmic rate for positive measures]\label{cd:thm:rate}
There are absolute constants $c_0,c_1,c_2>0$ such that the following holds. Let $\mu,\mu'$ be positive measures of mass at most $M$, let $a>0$, and let $\eta=\|P_a*(\mu-\mu')\|_{L^1}$ be small enough that $L=\log(1/\eta)\ge e^2$. Then
\[
d_{\rm BL}(\mu,\mu')\le\frac{c_1Ma}{L-2\log L}+\frac{c_1(1+L)^{1/2}}{L^2}.
\]
Conversely, for every $\Omega\ge2$ there are positive measures $\mu,\mu'$ of mass $M$ with $\|P_a*(\mu-\mu')\|_{L^1}\le c_2M(1+a)^{1/2}e^{-a(\Omega-1)}$ and $d_{\rm BL}(\mu,\mu')\ge c_0M/\Omega$. So positivity does not improve the rate $a/\log(1/\eta)$ (proved).
\end{theorem}

The first term dominates, and the bound has the order $Ma/L$, when $Ma\gtrsim L^{-1/2}$.

\begin{proof}
Fix a smooth even $\Lambda$ with $0\le\Lambda\le1$, $\Lambda=1$ on $[-\tfrac12,\tfrac12]$ and support in $[-1,1]$, and let $k_\Omega$ have Fourier transform $\Lambda(\omega/\Omega)$, so $k_\Omega(t)=\Omega k(\Omega t)$ with $k$ a Schwartz function and $\int k=1$. Let $\varphi$ be admissible and $\Delta=\mu-\mu'$. Since $\varphi$ is $1$ Lipschitz, $|\varphi-\varphi*k_\Omega|\le a_1/\Omega$ with $a_1=\int|s||k(s)|\,ds$, so $|\int(\varphi-\varphi*k_\Omega)\,d\Delta|\le2Ma_1/\Omega$. Let $\tilde k$ have Fourier transform $F(\omega)=\Lambda(\omega/\Omega)e^{a|\omega|}$, so $\tilde k*P_a=k_\Omega$ and, $k_\Omega$ and $P_a$ being even, $\int(\varphi*k_\Omega)\,d\Delta=\int(\varphi*\tilde k)\,(P_a*\Delta)\,dy$, at most $\|\tilde k\|_{L^1}\eta$. Carlson's inequality $\|f\|_{L^1}\le a_2\|f\|_{L^2}^{1/2}\|tf\|_{L^2}^{1/2}$ and Plancherel give $\|\tilde k\|_{L^1}\le a_3\|F\|_{L^2}^{1/2}\|F'\|_{L^2}^{1/2}$. On the support $|F|\le e^{a\Omega}$ and $|F'|\le(\|\Lambda'\|_\infty/\Omega+a)e^{a\Omega}$, so $\|\tilde k\|_{L^1}\le a_4e^{a\Omega}(1+a\Omega)^{1/2}$. Therefore $|\int\varphi\,d\Delta|\le2Ma_1/\Omega+a_4e^{a\Omega}(1+a\Omega)^{1/2}\eta$. Choose $a\Omega=L-2\log L$, so that $e^{a\Omega}\eta=L^{-2}$.

For the converse take the Fej\'er density $w(y)=\frac{2M}\pi\bigl(\frac{\sin(y/2)}y\bigr)^2$, of mass $M$, whose Fourier transform is a triangle supported in $[-1,1]$ with height $M$. Put $\mu=w(1+\cos\Omega y)\,dy$ and $\mu'=w\,dy$. Both are positive and both have mass $M$, because $\hat w(\pm\Omega)=0$. The difference $w\cos\Omega y$ has Fourier transform $\tfrac12(\hat w(\omega-\Omega)+\hat w(\omega+\Omega))$, supported in $\Omega-1\le|\omega|\le\Omega+1$, where $e^{-a|\omega|}\le e^{-a(\Omega-1)}$. The transform of $P_a*(w\cos\Omega y)$ and its derivative are therefore bounded by $Me^{-a(\Omega-1)}$ and $M(1+a)e^{-a(\Omega-1)}$ on a set of length $4$, and Carlson's inequality gives the $L^1$ bound. The test function $\varphi=\Omega^{-1}\cos\Omega y$ is admissible, and $\int\varphi\,w\cos\Omega y\,dy=\frac1{2\Omega}\int w(1+\cos2\Omega y)\,dy=\frac M{2\Omega}$, since $\hat w(2\Omega)=0$.
\end{proof}

The windowed problem with an estimated far field, and measures of infinite mass such as the zero counting measure, remain open.


\begin{proposition}[saturation of positive kernels]\label{cd:prop:sat}
Let $R\ge0$ be even with $\int R=1$ and $\widehat R$ supported in $[-1,1]$, and put $R_\Omega(x)=\Omega R(\Omega x)$. Then $1-\widehat R(u)\ge c_Ru^2$ for $|u|\le1$ with $c_R>0$. Consequently, for every $\mu\in L^2$ with $\mu\ne0$ there are $c>0$ and $\Omega_0$ with
\[
\|R_\Omega*\mu-\mu\|_{L^2}\ge c\,\Omega^{-2}\qquad(\Omega\ge\Omega_0).
\]
If the bandwidth is limited to $\Omega\le a^{-1}\log(1/\delta)$, the level at which noise of size $\delta$ amplified by $e^{a\Omega}$ reaches one, the bias of any such positive linear method is at least of order $(a/\log(1/\delta))^2$. Status: proved.
\end{proposition}

\begin{proof}
$1-\widehat R(u)=\int(1-\cos ux)R(x)\,dx\ge\frac{2u^2}{\pi^2}\int_{|x|\le\pi}x^2R(x)\,dx$ for $|u|\le1$, using $1-\cos\theta\ge2\theta^2/\pi^2$ on $|\theta|\le\pi$, and the last integral is positive because $R$ is band limited, hence real analytic, so not zero almost everywhere on $[-\pi,\pi]$. By Plancherel, $\|R_\Omega*\mu-\mu\|_2^2=\frac1{2\pi}\int|1-\widehat R(\omega/\Omega)|^2|\widehat\mu(\omega)|^2d\omega\ge c_R^2\Omega^{-4}\frac1{2\pi}\int_{|\omega|\le\Omega_0}\omega^4|\widehat\mu(\omega)|^2d\omega$ for $\Omega\ge\Omega_0$, and $\Omega_0$ is chosen so that the last integral is positive.
\end{proof}

This is the Korovkin saturation phenomenon (\cite[Chapter~9]{DeVoreLorentz93}): nonnegative convolution operators cannot converge faster than $\Omega^{-2}$, the Fej\'er kernel reaches only $\Omega^{-1}$ and the Jackson kernel reaches $\Omega^{-2}$. Building positivity into the linear operator therefore turns the H\"older rate of Theorem~\ref{cd:thm:holder} into a logarithmic one.

\begin{proposition}[the positive part is a pointwise contraction]\label{cd:prop:pospart}
If $\mu\ge0$ is a function and $f$ is any real estimate of it, then $|f_+-\mu|\le|f-\mu|$ pointwise, so $f_+\ge0$ has every $L^p$ error bound of $f$. Status: proved.
\end{proposition}

\begin{proof}
Where $f\ge0$ nothing changes. Where $f<0\le\mu$, $|0-\mu|\le|f-\mu|$.
\end{proof}

Positivity can therefore be imposed after a fast linear inversion, by taking the positive part, without loss of accuracy; the same holds for the projection in $L^2$ onto any closed convex set that contains the truth. Both propositions are from an unrefereed note of Polson (November 2026), and both are classical.

At the boundary the inverse loses all but a logarithm for a source in a Sobolev ball (Theorem~\ref{cd:thm:log}), which the zero measure is not.

\section{Statistics}\label{cd:sec:statistics}

\subsection{Supersmooth noise}

In density deconvolution one observes $Y_i=X_i+E_i$, $i\le n$, with $E_i$ of known law, and estimates the density $f$ of $X_i$; the difficulty is set by the decay of the characteristic function $\phi_E$ of the noise. Cauchy noise has $\phi_E(\omega)=e^{-a|\omega|}$ and is supersmooth of order one in the classification of Fan \cite{Fan91}, who showed that for supersmooth noise the minimax rate is logarithmic in $n$ (cited). Carroll and Hall \cite{CarrollHall88} found this for Gaussian noise before Fan. Stefanski and Carroll \cite{StefanskiCarroll90} introduced the deconvoluting kernel estimators, which attain these rates \cite{CarrollHall88,Fan91}; the estimator has a Fourier cutoff at a frequency of order $\log n$. Heuristically, the logarithm is the statistical form of Theorem~\ref{cd:thm:log} with $\varepsilon=n^{-1/2}$.

The rates improve when the target is smooth in the strong sense. Butucea and Tsybakov \cite{ButuceaTsybakov08a,ButuceaTsybakov08b} constructed sharp asymptotic minimax estimators when the target is itself supersmooth, with rates much faster than the logarithmic ones (cited). Cauchy noise with a target analytic in a strip is simpler, a bound of the form of Proposition~\ref{cd:prop:finite}, stated here without attribution.

\begin{proposition}[a strip]\label{cd:prop:strip}
Let $E_i$ be Cauchy noise of width $a$ and let $f$ satisfy $\int e^{2w|\omega|}|\widehat f(\omega)|^2d\omega\le2\pi M^2$ with $w>0$. The spectral cutoff estimator $\widehat f_n(x)=(2\pi)^{-1}\int_{|\omega|\le\Omega}e^{a|\omega|}\widehat\phi_n(\omega)e^{i\omega x}d\omega$, where $\widehat\phi_n(\omega)=n^{-1}\sum_{j\le n}e^{-i\omega Y_j}$ is the empirical transform of $Y_1,\dots,Y_n$ in the convention of Section~\ref{cd:sec:operator}, with $\Omega=\log n/(2(a+w))$, has root mean integrated squared error $O(n^{-w/(2(a+w))})$. This is an upper bound.
\end{proposition}

\begin{proof}
The transform of the law of $Y$ is $\widehat f(\omega)e^{-a|\omega|}$ and $\mathbb E\lvert\widehat\phi_n(\omega)-\widehat f(\omega)e^{-a|\omega|}\rvert^2\le1/n$. By Plancherel, the integrated variance is at most $(2\pi n)^{-1}\int_{|\omega|\le\Omega}e^{2a|\omega|}d\omega\le e^{2a\Omega}/(2\pi an)$, and the integrated squared bias is $(2\pi)^{-1}\int_{|\omega|>\Omega}|\widehat f|^2\le M^2e^{-2w\Omega}$. With $e^{2(a+w)\Omega}=n$ both are of order $n^{-w/(a+w)}$.
\end{proof}

A source of Lorentzians of width $w$ (the test source of Table~\ref{cd:tab:slopes}) is analytic in every strip of width less than $w$, and the exponent $(b+w)/(a+w)$ of Proposition~\ref{cd:prop:finite} is positive even at $b=0$. Meister's book \cite{Meister09} collects the theory.

\subsection{Bayesian posteriors}\label{cd:sec:bayes}

For linear inverse problems with Gaussian priors, Knapik, van der Vaart and van Zanten \cite{KnapikVZ11} computed the posterior contraction rates in terms of the decay of the singular values, for polynomially decaying singular values only, the mildly ill posed case (cited). Cauchy deconvolution is severely ill posed: on a periodic domain its singular values are $e^{-a\lvert k\rvert}$ (convolution with $P_a$ on $L^2(\R)$ is not compact, being bounded below on the functions with spectrum in $[-1,1]$, so a periodic or windowed setting is needed). The severely ill posed case was treated by the same authors for the recovery of the initial condition of the heat equation \cite{KnapikVZ13}, the Gaussian analogue of the backward Cauchy flow. For mixtures, where the unknown is a mixing measure seen through a known kernel, Nguyen \cite{Nguyen13} gave rates in Wasserstein distance for mixing measures with compact support: polynomial for ordinary smooth kernels and logarithmic for supersmooth ones, of order $(\log n)^{-1/\beta}$ for a kernel of supersmooth order $\beta$, $(\log n)^{-1/2}$ for the Gaussian kernel; for the Cauchy kernel, $\beta=1$, the hypotheses of the paper are not checked here. Gao and van der Vaart \cite{GaoVdV16} treated Dirichlet process mixtures of Laplace kernels (the noise is Laplace, not Cauchy). No posterior theorem for Cauchy deconvolution is cited. By analogy the Bayesian picture is the deterministic one with $\varepsilon$ the noise level; this is a heuristic, since prior mass, testing and support conditions enter.

\subsection{Positivity in the algorithms}

Spectral cutoff, Tikhonov regularisation and Landweber iteration are linear and do not: they return signed estimates even when a positive source exists. The positivity preserving methods are nonlinear: nonnegative least squares (Section~\ref{sec:ch17:nnls}), maximum entropy, and the EM algorithm of Richardson and Lucy \cite{Richardson72,Lucy74}, studied for positive linear inverse problems by Vardi and Lee \cite{VardiLee93} (Section~\ref{lo:sec:bayes}). Chapters~\ref{chap:lo} and~\ref{chap:cx} use EM, and its accelerated form in the logarithmic domain (Section~\ref{cx:sec:accel}), for this reason. When no positive source exists, which by Proposition~\ref{cd:prop:critical} is the case beyond the critical width, these methods still return positive estimates, each its own projection (least squares for nonnegative least squares, Kullback Leibler for EM), and the residual, compared with the data error, is what detects the failure (Proposition~\ref{lo:prop:emres}): positivity of the method is not evidence of positivity of the source.

\section{Atomic sources}\label{cd:sec:atomic}

\subsection{From logarithmic to polynomial}

The rates above are for smoothness classes. For a finite sum of atoms, $\mu=\sum_kc_k\delta_{t_k}$, the problem is different in kind. The data are a sum of Lorentzians, $\sum_kc_kP_a(x-t_k)$, whose transform is the sum of exponentials $\sum_kc_ke^{-a|\omega|-i\omega t_k}$, and after a taper, as in Theorem~\ref{es:thm:taper}, the samples on an arithmetic progression of frequencies are a sum of damped complex exponentials. Recovering their poles is Prony's problem, solved stably by Hankel matrix methods: the matrix pencil of Hua and Sarkar \cite{HuaSarkar90}, and the Hankel singular value decomposition used in nuclear magnetic resonance spectroscopy since Barkhuijsen, de Beer and van Ormondt \cite{Barkhuijsen87}, where Lorentzian line shapes and their damped exponential transforms are the daily object. For a fixed number of separated atoms with amplitudes bounded below and no aliasing, the error is Lipschitz in the data error, with a constant that grows polynomially as atoms approach each other, except near merging atoms, where it is the square root; this is super resolution, studied by Donoho \cite{Donoho92}, Cand\`es and Fernandez Granda \cite{CandesFG14}, and, for positive measures without a separation condition, Schiebinger, Robeva and Recht \cite{SchiebingerRR18}, with sharp constants for close atoms by Li and Liao \cite{LiLiao21} (cited). Chapter~\ref{chap:es} compares EM on a grid with the pencil on the zeros of $\zeta$.

\subsection{Merging atoms}

The polynomial rate has one exception that matters for $\zeta$. If two atoms approach each other symmetrically, the data depend on their separation only through its square, the Fisher information for the separation vanishes, and the rate for the separation is the square root of the data error (with more merging atoms the exponent can be smaller). It is the singular rate of finite mixtures (Chen \cite{Chen95}; Nguyen \cite{Nguyen13}) in its simplest case, and Chapter~\ref{chap:es} met it in the form forced by the functional equation: a zero off the line comes with its partner at the same ordinate, their widths at level $\theta$ are $\theta-\tfrac12\mp\delta$, the two poles merge as $\delta\to0$, and a data error $\tau$ determines $\delta$ only to order $\tau^{1/2}$ (Lemma~\ref{es:lem:even} and Theorem~\ref{es:thm:singular}).

The square root rate describes the error of a fit. A second question is what the data cannot exclude, and a bound in that direction follows from matching moments (an unrefereed note of Polson, November 2026): positive measures with the same first four moments are nearly indistinguishable after blurring.

\begin{proposition}[A cluster that imitates a pair]\label{cd:prop:cluster}
Let $a>0$, $\pi_\delta=\delta_{-\delta/2}+\delta_{\delta/2}$ and $\kappa_\delta=\tfrac14\delta_{-\delta}+\tfrac32\delta_0+\tfrac14\delta_\delta$. Both are positive of mass $2$ with the same moments of orders $0$ to $3$, and
\[
\frac{\|P_a*(\pi_\delta-\kappa_\delta)\|_\infty}{2P_a(0)}\le\frac5{16}\Bigl(\frac\delta a\Bigr)^4.
\]
Hence if the blurred data of a pair are known to relative accuracy $\eta$, that is to $\eta$ times the peak height $2P_a(0)$, they cannot exclude the three atom cluster when $\delta\le a(16\eta/5)^{1/4}$; for $a=\tfrac12$ and $\eta=2\times10^{-4}$ this is about $0.08$ (proved, and computed for the number).
\end{proposition}

\begin{proof}
The moments of orders $0,1,2,3$ of both measures are $2,0,\delta^2/2,0$. Since $P_a(y)=\frac1\pi\operatorname{Im}(y-ia)^{-1}$, $|P_a^{(k)}|\le k!/(\pi a^{k+1})$. Expand $P_a(y-t)$ in $t$ to order three, with remainder at most $t^4/(\pi a^5)$. The polynomial parts cancel for the two measures, and the remainders sum to at most $\bigl(2(\tfrac\delta2)^4+\tfrac12\delta^4\bigr)/(\pi a^5)=\frac{5\delta^4}{8\pi a^5}$. Divide by $2P_a(0)=2/(\pi a)$.
\end{proof}

This is a bound on what can be identified. Below a gap of order $a\eta^{1/4}$, polynomial in the data error and unlike the worst case of Theorem~\ref{cd:thm:rate}, no method can tell the pair from a cluster. It says nothing about gaps above that order: that larger gaps are identifiable is Problem~\ref{cd:prob:cluster}. The order $(\delta/a)^4$ is plausible, since a nearby positive competitor can match only the moments below the order at which the two atom measure becomes determinate (heuristic).

\subsection{Analytic continuation in physics}

A related problem is solved daily in many body physics, where a spectral function $A(\omega)\ge0$ is to be recovered from its Cauchy transform $G(z)=\int A(\omega)(z-\omega)^{-1}d\omega$ known at points away from the real axis, the Matsubara frequencies. Its imaginary part at height $y$ is $-\pi(P_y*A)$, but the observation operator differs from the one above: the data sample the Poisson integral at $x=0$ and several heights, not along a line, so uniqueness and stability do not transfer (finite data are not unique). It is severely ill posed, unconstrained methods produce negative spectral weight, and the standard tools are maximum entropy (Jarrell and Gubernatis \cite{JarrellGubernatis96}) and, more recently, Nevanlinna interpolation (Fei, Yeh and Gull \cite{FeiYehGull21}), which builds positivity in by parametrising the Pick functions of Theorem~\ref{cd:thm:herglotz} directly. Section~\ref{cx:sec:vonneumann} gives the identity that links the two: the Thorin density is $u_\theta(v)=\pi^{-1}\operatorname{Tr}\bigl[a((H-v)^2+a^2)^{-1}\bigr]$ with $a=\theta-\tfrac12$, for a hypothetical self adjoint $H$ whose spectrum is the set of ordinates (so under RH); the resolvent of $H$ itself is not trace class. 

\section{What is known and what is not}\label{cd:sec:known}

Two features are specific to the book's use of the operator. The first is the systematic use of the whole family of partial deconvolutions, the Thorin family $\mathfrak U_\theta=P_{\theta-\theta'}*\mathfrak U_{\theta'}$, with positivity at each level, and its critical width $c^*(u_1)=1-\sigma_{\max}$ (Theorem~\ref{tf:thm:deconv}(a)). The second is the identification of the full deconvolution of the level one density with Weil's explicit formula (Theorem~\ref{dc:thm:weil} and Chapter~\ref{chap:es}), which is classical in substance.

\begin{problem}\label{cd:prob:width}
Estimating the critical width. Given noisy data $u+\text{noise}$ with $u\ge0$ and $\int u(x)(1+x^2)^{-1}dx<\infty$, estimate $c^*(u)$ of Proposition~\ref{cd:prop:critical}. For atomic sources with positive weights it is the smallest width among the poles, estimable by the pencil on a class with bounded order, weights bounded below, separation and no aliasing; without weights bounded below there is no uniform rate, since $c^*$ is discontinuous ($(1-\delta)P_A+\delta P_r$ has $c^*=r<A$ for every $\delta>0$). For general sources it is a boundary of a convex cone in an infinite dimensional space. Status: open.
\end{problem}

\begin{problem}\label{cd:prob:rates}
Positivity and rates. Positivity does not improve the logarithmic minimax rate over smoothness classes of densities in the worst case (the lower bounds of Fan \cite{Fan91} are over such classes), and sparsity improves it to a polynomial one. Imposing a prescribed number of sign changes of a real Fourier transform does not help: two positive densities whose transforms are $e^{-\lvert\omega\rvert}$ and that function plus two nonnegative bumps near $\pm\Omega$ both have no sign change, and keep the logarithmic obstruction. Smoothing and a mass bound give a polynomial exponent (Proposition~\ref{cd:prop:finite} at $b=0$), which positivity does not. For positive measures of mass at most $M$ the logarithmic rate is sharp (Theorem~\ref{cd:thm:rate}). Open: is the Thorin family in a class on which positivity alone gives a polynomial rate?
\end{problem}

\begin{problem}\label{cd:prob:three}
Three lines with positivity on a window. Theorem~\ref{cd:thm:holder} needs an a priori bound at a level below the target, and Corollary~\ref{cd:cor:positive} removes it for data in $L^2$ on the whole line. For data known only on a window, the mass in a window of a positive source is controlled by the data alone, as in the a posteriori bound of Proposition~\ref{cx:prop:apost}. Does positivity give a H\"older bound between levels from data on a window, without any a priori bound below, with an exponent determined by the data? Status: open.
\end{problem}

\begin{problem}\label{cd:prob:cluster}
Close pairs. Prove that the bound of Proposition~\ref{cd:prop:cluster} is sharp: any positive measure of mass $2$ near $\pi_\delta$ whose blurred version fits that of $\pi_\delta$ to relative accuracy $\eta$ lies within $C(a/\delta)^4\eta$ of $\pi_\delta$ in the bounded Lipschitz distance. Status: open.
\end{problem}

\subsection*{Notes}

The bibliography entries that are new for this chapter have been checked against CrossRef or OpenAlex, except \cite{Hadamard23}, and the statements about the contents of four of the papers (Knapik et al. 2011 and 2013, Butucea and Tsybakov, Fan, Gao and van der Vaart) are those of a web check made for this chapter and about the others those of the reviews, not a reading of the papers here (in particular the supersmooth order $\beta$ of Nguyen \cite{Nguyen13}, and the form of the rates of \cite{KnapikVZ13}, which is not stated).
